\section{The \instname\ data reduction flow.}
\label{sec:reduction_flow}


The overall data flow of the \instname\ pipeline is displayed in Figure \ref{fig:reduction_cascade}.

 
\begin{figure}[h!]
\centering
\includegraphics[width=0.9\textwidth]{../instrument_tutorial/figures/reduction_cascade.jpg}
\caption{The data reduction cascade of the \instname\ workflow.}
\label{fig:reduction_cascade}
\end{figure}



The reduction cascade is organized in tasks, which represent
well-defined steps in the process. Tasks can be grouped inside
sub-workflows.  Each task runs a recipe; the detailed description of
the algorithms, input, outputs and recipe parameters used in each
recipe are available in the pipeline manual. Here, we present only the
description of most important features.

The \wkfname\  \edps\ workflow is designed to execute the tasks that
deliver the final reduced data cube for each dataset. It can be either
the product of a single exposure, or the combination of multiple
exposures. Only calibrations needed by the selected the scientific
exposures are processed.

It is possible to set \edps\ to perform the data reduction until a
certain step of the reduction chain (e.g. to reduce only standar
stars, or only flat fields).  This is done by specifying the desired
tasks in the field {\bf Select reduction target} of the {\bf Raw Data}
tab.

The reduction steps of the \wkfname\ workflow are listed below. Before
starting the reduction, the parameters of the recipes associated to
each task can be configured by pressing the button

\includegraphics[width=0.8cm]{../common/figures/configure_dataset.jpg}
close to each dataset configuration.
See for more info on the configuration editor \ref{sec:configuration_editor}



\subsection{Imaging, Coronography, Sparse aperture masking}
\begin{itemize}
\item Data repacking and dark/sky subtraction 
({\tt visir\_util\_repack}). 
The input data is first split – the science images and, if present, 
the photometric standards are sent to be decoded, this step 
transforms the input FITS files into a format (still FITS) that 
makes the subsequent processing easier. The chop/nod background 
subtraction is applied.
\item Detecting objects and calculating shifts
({\tt visir\_util\_detect\_shift}). This is critical for the 
subsequent image alignment. A sky background 
frame is created during this step, after applying a kappa-sigma 
clipping algorithm to the chop-nod corrected data, and finally the 
sky-subtracted images will be aligned with the appropriate offsets 
and co-added together to produce a final image.
\item Shifting and adding individual frames together 
({\tt visir\_util\_run\_swarp}).
\item Optional photometric calibration that is carried out only 
if the input Data Set contains a photometric standard. Only the 
final combined image is calibrated, so if the users want to study 
the variability of a source during the observations, they must 
split the data into separate Data Sets, process them separately, 
and they will be calibrated individually.
The stellar flux is corrected for the background, and the error 
budget includes the sky error. The aperture that maximizes the 
signal-to-noise ratio is used. The measured standard star flux 
is compared with the cataloged flux, in Jy. The final 
reported sensitivity is computed in units of mJy at 10-$\sigma$ 
level in 1\,hr (the conversion factor is the reciprocal of this 
value). If a conversion factor is derived for the Data Set, then 
the final fits are converted from ADU sec$^{-1}$ into physical 
units Jy.
For more details see the \instname\ pipeline manual.
\end{itemize}

\subsection{Spectroscopy modes}
\begin{itemize}
\item Data repacking ({\tt visir\_util\_repack}): the same as
the imaging case above.
\item Distortion Correction ({\tt visir\_util\_undistort}) is 
based on analytical parametrization. The parametrization was 
derived from an instrument model. The correction is applied 
directly by interpolating the distortion corrected pixel value 
from the source pixels. The pixels that are marked as bad are 
ignored during the interpolation.
\item Wavelength Calibration and Spectral Extraction
({\tt visir\_util\_detect\_shift},\\ {\tt visir\_old\_spc\_wcal},
{\tt visir\_old\_spc\_wcal\_ech}, {\tt visir\_old\_spc\_obs},\\ 
{\tt visir\_old\_spc\_obs\_ech}), based on a polynomial that 
describes the dispersion relation. The \instname\ physical 
model -- represented by a first degree polynomial -- is used 
as an initial guess. Experience has shown that only a linear 
shift is required and this is achieved by cross-correlating a 
sky spectrum derived from the data and a model of the atmospheric 
emission. The derived offset is then applied to the model.
By default the pipeline extracts a 1D spectrum of the brightest 
source, applying the optimal extraction method. 
%However, the 
%users have options to define on each image multiple (up to 10) 
%extraction apertures and to use different methods - optimal or 
%aperture extraction - to obtain a 1D spectrum. Therefore, the 
%pipeline allows simultaneous (in a single run) 1D spectra 
%extraction from different objects that fall onto the slit or 
%from different parts of the same object (e.g. core, jet, disk 
%of a galaxy, etc.).
\item Photometric Calibration ({\tt visir\_old\_spc\_phot}, 
{\tt visir\_old\_spc\_phot\_ech}) is carried out by\\ comparing 
a flux model (available in the static calibration database) of 
the standard star with the observed flux. The final reported 
sensitivity is computed in units of mJy at the 10-$\sigma$ level 
in 1\,hr (the conversion factor is the reciprocal of this value). 
The normalization per unit time is obtained calculating the 
total integration time as follows: 
{\tt DIT$\times$NDIT$\times$NFILES$\times$NCHOP$\times$2}. The 
factor of two is due to the half-cycle chopping. For more 
details, see the \instname\ pipeline manual.
\end{itemize}

For further details see \instname\ Pipeline User Manual, Sec.\,8.
Note that the \instname\ pipeline cascade omits two steps that 
the users might expect to see:
\begin{itemize}
\item separate dark subtraction step -- because the dark signal 
is subtracted together with the sky emission;
\item flat-fielding, -- because the detector is too unstable, 
and the flat-fielding procedure would only add noise.
\end{itemize}
